% ═══════════════════════════════════════════════════════════════════════
\section{Synthèse}
% ═══════════════════════════════════════════════════════════════════════

\begin{frame}{Ce qu'il faut retenir}
  \begin{itemize}
    \item Une transformation = une \textbf{bijection} du plan ; toutes
      les transformations \textbf{affines} tiennent dans
      $w(P) = AP + \vec t$.
    \item Une chaîne de sous-groupes emboîtés
      $\Iso^+ \subset \Iso \subset \Sim \subset \Aff \subset \Transf$ :
      plus le groupe est petit, plus il conserve de propriétés (Klein).
    \item Le carré unité et ses colonnes-images : la façon la plus
      rapide de \emph{voir} ce que fait une matrice.
    \item Les coordonnées homogènes rendent l'affine linéaire, une
      dimension plus haut --- toute l'infographie repose là-dessus.
    \item Trois points déterminent une application affine ; deux
      points, une similitude ; un point, une translation.
    \item Toute isométrie est translation, rotation, réflexion ou
      réflexion glissée --- et la réflexion engendre tout.
  \end{itemize}
\end{frame}

\begin{frame}{Figure 1.3, une dernière fois}
  \FigBox[0.62\textheight]{%
    \begin{tikzpicture}[scale=1.05, every node/.style={font=\small}]
      \foreach \i/\lab/\prop in {%
         0/{$\Transf$ : toutes les transformations}/{aucun invariant métrique},
         1/{$\Aff$ : transformations affines}/{alignement, parallélisme},
         2/{$\Sim$ : similitudes}/{$+$ angles, rapports de longueurs},
         3/{$\Iso$ : isométries}/{$+$ longueurs, aires},
         4/{$\Iso^{+}$ : déplacements}/{$+$ orientation}}
        {\pgfmathsetmacro{\w}{9.4-1.35*\i}
         \pgfmathsetmacro{\y}{-0.95*\i}
         \draw[rounded corners=2pt, fill=steelblue4!\the\numexpr8+10*\i\relax,
               draw=sectioncolor!60]
           ({-\w/2},{\y-0.36}) rectangle ({\w/2},{\y+0.36});
         \node at (0,\y) {\lab};
         \node[anchor=west, font=\scriptsize, gray!80!black]
           at ({\w/2+0.25},\y) {\prop};}
    \end{tikzpicture}}
  \small Chaque étage de cette pyramide a été démontré (ou admis)
  aujourd'hui.
\end{frame}

\begin{frame}{Et maintenant}
  \begin{itemize}
    \item Les \textbf{démonstrations} laissées « au tableau »
      aujourd'hui sont toutes dans les notes de cours, avec tous les
      détails.
    \item Les \textbf{exercices} du chapitre (avec solutions
      complètes) couvrent chacun des résultats vus : à travailler
      avant le prochain cours.
    \item L'\textbf{annexe Python} des notes reprend en code toutes
      les idées centrales (le bulldozer, la reconnaissance d'une
      similitude, la décomposition d'une isométrie).
  \end{itemize}
\end{frame}
